% ============================================================
%  TEMPLATE BEAMER --- Tema "Dark Accent"
% ============================================================
%  Este arquivo concentra TODO o preâmbulo compartilhado entre
%  as aulas (classe do documento, pacotes, cores, fontes,
%  cabeçalho/rodapé dos slides e comandos auxiliares).
%
%  Como usar em um novo arquivo de aula:
%
%      % ============================================================
%  TEMPLATE BEAMER --- Tema "Dark Accent"
% ============================================================
%  Este arquivo concentra TODO o preâmbulo compartilhado entre
%  as aulas (classe do documento, pacotes, cores, fontes,
%  cabeçalho/rodapé dos slides e comandos auxiliares).
%
%  Como usar em um novo arquivo de aula:
%
%      % ============================================================
%  TEMPLATE BEAMER --- Tema "Dark Accent"
% ============================================================
%  Este arquivo concentra TODO o preâmbulo compartilhado entre
%  as aulas (classe do documento, pacotes, cores, fontes,
%  cabeçalho/rodapé dos slides e comandos auxiliares).
%
%  Como usar em um novo arquivo de aula:
%
%      % ============================================================
%  TEMPLATE BEAMER --- Tema "Dark Accent"
% ============================================================
%  Este arquivo concentra TODO o preâmbulo compartilhado entre
%  as aulas (classe do documento, pacotes, cores, fontes,
%  cabeçalho/rodapé dos slides e comandos auxiliares).
%
%  Como usar em um novo arquivo de aula:
%
%      \input{templates/beamer.tex}
%
%      \title{Título da Aula}
%      \author
%
%      \begin{document}
%        ... seus \begin{frame} ... \end{frame} ...
%      \end{document}
%
%  IMPORTANTE: como este arquivo contém o \documentclass, ele
%  precisa ser a PRIMEIRA linha do .tex da aula (nenhum outro
%  comando pode vir antes do \input acima).
%
%  Compilar sempre a partir da pasta latex/ (para o caminho
%  relativo "templates/beamer.tex" ser encontrado):
%      cd latex && pdflatex aula01.tex
% ============================================================

% !TEX root = ../aula01.tex

\documentclass[11pt, aspectratio=169, xcolor=table]{beamer}

% ===================== PACOTES =====================
\usepackage[utf8]{inputenc}
\usepackage[T1]{fontenc}
\usepackage{lmodern}
\usepackage{newtxtext}
\usepackage{newtxmath}
\usepackage{tikz}
\usepackage{amsmath, amssymb}
\usepackage{booktabs}
\usepackage{bookmark}

% Atalho para "d" reto de diferenciais: dx, dy -> \dd x, \dd y
\newcommand{\dd}{\mathrm{d}}

% ===================== REMOVE ELEMENTOS PADRÃO DO BEAMER =====================
\setbeamertemplate{navigation symbols}{}

% ===================== CORES =====================
\definecolor{bgcolor}{RGB}{20, 20, 24}
\definecolor{accent}{RGB}{214, 165, 113}
\definecolor{dotoff}{RGB}{70, 70, 76}

\setbeamercolor{background canvas}{bg=bgcolor}
\setbeamercolor{normal text}{fg=white!92!black}
\setbeamercolor{title}{fg=white}
\setbeamercolor{frametitle}{fg=accent}
\setbeamercolor{footline}{fg=gray!60}
\setbeamercolor{itemize item}{fg=accent}
\setbeamercolor{itemize subitem}{fg=accent!70!white}
\setbeamercolor{progress}{bg=bgcolor}
\setbeamercolor{footnum}{bg=bgcolor}
\setbeamertemplate{itemize item}{\textbullet}

\setbeamerfont{frametitle}{series=\bfseries}
\setbeamerfont{footline}{size=\tiny}

\hypersetup{colorlinks=true, urlcolor=accent, linkcolor=accent}

% ===================== TÍTULO DO SLIDE =====================
\setbeamertemplate{frametitle}{
  \vspace{0.5cm}
  \hspace{0.4cm}\usebeamerfont{frametitle}\usebeamercolor[fg]{frametitle}\insertframetitle\\
  \vspace{0.12cm}
  \hspace{0.4cm}\textcolor{accent}{\rule{2.5cm}{0.6pt}}
}

% ===================== RODAPÉ =====================
\setbeamertemplate{footline}{%
  \begin{beamercolorbox}[wd=\paperwidth, ht=0.09cm, dp=0pt]{progress}
    \begin{tikzpicture}
      \fill[dotoff] (0,0) rectangle (16,0.09);
      \ifnum\inserttotalframenumber>0
        \pgfmathsetmacro{\prog}{16*\insertframenumber/\inserttotalframenumber}
        \fill[accent] (0,0) rectangle (\prog,0.09);
      \fi
    \end{tikzpicture}
  \end{beamercolorbox}%
  \vskip 2pt
  \begin{beamercolorbox}[wd=\paperwidth, ht=2.4ex, dp=1.2ex, right]{footnum}
    \usebeamerfont{footline}\color{gray!60}\insertframenumber\,/\,\inserttotalframenumber\hspace{0.8em}
  \end{beamercolorbox}%
}

\setlength{\parindent}{0pt}
\setlength{\parskip}{6pt}

% ===================== CUSTOMIZAÇÃO DE BLOCOS =====================
\setbeamertemplate{blocks}[rounded][shadow=false]
\setbeamercolor{block title}{bg=dotoff, fg=accent}
\setbeamercolor{block body}{bg=bgcolor!90!white, fg=white!92!black}

\setbeamercolor{block title alerted}{bg=accent!60!black, fg=white}
\setbeamercolor{block body alerted}{bg=bgcolor!95!white, fg=white!92!black}

% ===================== CAIXA DE DESTAQUE =====================
% Uso: \notebox{texto da observação}
\newcommand{\notebox}[1]{%
  \begin{center}
  \begin{tabular}{@{}l@{\hspace{0.35cm}}p{0.82\textwidth}@{}}
    \textcolor{accent}{\rule{2pt}{0.55cm}} & \itshape\small\color{white!88!black} #1
  \end{tabular}
  \end{center}%
}

\usefonttheme[onlymath]{serif}
%
%      \title{Título da Aula}
%      \author
%
%      \begin{document}
%        ... seus \begin{frame} ... \end{frame} ...
%      \end{document}
%
%  IMPORTANTE: como este arquivo contém o \documentclass, ele
%  precisa ser a PRIMEIRA linha do .tex da aula (nenhum outro
%  comando pode vir antes do \input acima).
%
%  Compilar sempre a partir da pasta latex/ (para o caminho
%  relativo "templates/beamer.tex" ser encontrado):
%      cd latex && pdflatex aula01.tex
% ============================================================

% !TEX root = ../aula01.tex

\documentclass[11pt, aspectratio=169, xcolor=table]{beamer}

% ===================== PACOTES =====================
\usepackage[utf8]{inputenc}
\usepackage[T1]{fontenc}
\usepackage{lmodern}
\usepackage{newtxtext}
\usepackage{newtxmath}
\usepackage{tikz}
\usepackage{amsmath, amssymb}
\usepackage{booktabs}
\usepackage{bookmark}

% Atalho para "d" reto de diferenciais: dx, dy -> \dd x, \dd y
\newcommand{\dd}{\mathrm{d}}

% ===================== REMOVE ELEMENTOS PADRÃO DO BEAMER =====================
\setbeamertemplate{navigation symbols}{}

% ===================== CORES =====================
\definecolor{bgcolor}{RGB}{20, 20, 24}
\definecolor{accent}{RGB}{214, 165, 113}
\definecolor{dotoff}{RGB}{70, 70, 76}

\setbeamercolor{background canvas}{bg=bgcolor}
\setbeamercolor{normal text}{fg=white!92!black}
\setbeamercolor{title}{fg=white}
\setbeamercolor{frametitle}{fg=accent}
\setbeamercolor{footline}{fg=gray!60}
\setbeamercolor{itemize item}{fg=accent}
\setbeamercolor{itemize subitem}{fg=accent!70!white}
\setbeamercolor{progress}{bg=bgcolor}
\setbeamercolor{footnum}{bg=bgcolor}
\setbeamertemplate{itemize item}{\textbullet}

\setbeamerfont{frametitle}{series=\bfseries}
\setbeamerfont{footline}{size=\tiny}

\hypersetup{colorlinks=true, urlcolor=accent, linkcolor=accent}

% ===================== TÍTULO DO SLIDE =====================
\setbeamertemplate{frametitle}{
  \vspace{0.5cm}
  \hspace{0.4cm}\usebeamerfont{frametitle}\usebeamercolor[fg]{frametitle}\insertframetitle\\
  \vspace{0.12cm}
  \hspace{0.4cm}\textcolor{accent}{\rule{2.5cm}{0.6pt}}
}

% ===================== RODAPÉ =====================
\setbeamertemplate{footline}{%
  \begin{beamercolorbox}[wd=\paperwidth, ht=0.09cm, dp=0pt]{progress}
    \begin{tikzpicture}
      \fill[dotoff] (0,0) rectangle (16,0.09);
      \ifnum\inserttotalframenumber>0
        \pgfmathsetmacro{\prog}{16*\insertframenumber/\inserttotalframenumber}
        \fill[accent] (0,0) rectangle (\prog,0.09);
      \fi
    \end{tikzpicture}
  \end{beamercolorbox}%
  \vskip 2pt
  \begin{beamercolorbox}[wd=\paperwidth, ht=2.4ex, dp=1.2ex, right]{footnum}
    \usebeamerfont{footline}\color{gray!60}\insertframenumber\,/\,\inserttotalframenumber\hspace{0.8em}
  \end{beamercolorbox}%
}

\setlength{\parindent}{0pt}
\setlength{\parskip}{6pt}

% ===================== CUSTOMIZAÇÃO DE BLOCOS =====================
\setbeamertemplate{blocks}[rounded][shadow=false]
\setbeamercolor{block title}{bg=dotoff, fg=accent}
\setbeamercolor{block body}{bg=bgcolor!90!white, fg=white!92!black}

\setbeamercolor{block title alerted}{bg=accent!60!black, fg=white}
\setbeamercolor{block body alerted}{bg=bgcolor!95!white, fg=white!92!black}

% ===================== CAIXA DE DESTAQUE =====================
% Uso: \notebox{texto da observação}
\newcommand{\notebox}[1]{%
  \begin{center}
  \begin{tabular}{@{}l@{\hspace{0.35cm}}p{0.82\textwidth}@{}}
    \textcolor{accent}{\rule{2pt}{0.55cm}} & \itshape\small\color{white!88!black} #1
  \end{tabular}
  \end{center}%
}

\usefonttheme[onlymath]{serif}
%
%      \title{Título da Aula}
%      \author
%
%      \begin{document}
%        ... seus \begin{frame} ... \end{frame} ...
%      \end{document}
%
%  IMPORTANTE: como este arquivo contém o \documentclass, ele
%  precisa ser a PRIMEIRA linha do .tex da aula (nenhum outro
%  comando pode vir antes do \input acima).
%
%  Compilar sempre a partir da pasta latex/ (para o caminho
%  relativo "templates/beamer.tex" ser encontrado):
%      cd latex && pdflatex aula01.tex
% ============================================================

% !TEX root = ../aula01.tex

\documentclass[11pt, aspectratio=169, xcolor=table]{beamer}

% ===================== PACOTES =====================
\usepackage[utf8]{inputenc}
\usepackage[T1]{fontenc}
\usepackage{lmodern}
\usepackage{newtxtext}
\usepackage{newtxmath}
\usepackage{tikz}
\usepackage{amsmath, amssymb}
\usepackage{booktabs}
\usepackage{bookmark}

% Atalho para "d" reto de diferenciais: dx, dy -> \dd x, \dd y
\newcommand{\dd}{\mathrm{d}}

% ===================== REMOVE ELEMENTOS PADRÃO DO BEAMER =====================
\setbeamertemplate{navigation symbols}{}

% ===================== CORES =====================
\definecolor{bgcolor}{RGB}{20, 20, 24}
\definecolor{accent}{RGB}{214, 165, 113}
\definecolor{dotoff}{RGB}{70, 70, 76}

\setbeamercolor{background canvas}{bg=bgcolor}
\setbeamercolor{normal text}{fg=white!92!black}
\setbeamercolor{title}{fg=white}
\setbeamercolor{frametitle}{fg=accent}
\setbeamercolor{footline}{fg=gray!60}
\setbeamercolor{itemize item}{fg=accent}
\setbeamercolor{itemize subitem}{fg=accent!70!white}
\setbeamercolor{progress}{bg=bgcolor}
\setbeamercolor{footnum}{bg=bgcolor}
\setbeamertemplate{itemize item}{\textbullet}

\setbeamerfont{frametitle}{series=\bfseries}
\setbeamerfont{footline}{size=\tiny}

\hypersetup{colorlinks=true, urlcolor=accent, linkcolor=accent}

% ===================== TÍTULO DO SLIDE =====================
\setbeamertemplate{frametitle}{
  \vspace{0.5cm}
  \hspace{0.4cm}\usebeamerfont{frametitle}\usebeamercolor[fg]{frametitle}\insertframetitle\\
  \vspace{0.12cm}
  \hspace{0.4cm}\textcolor{accent}{\rule{2.5cm}{0.6pt}}
}

% ===================== RODAPÉ =====================
\setbeamertemplate{footline}{%
  \begin{beamercolorbox}[wd=\paperwidth, ht=0.09cm, dp=0pt]{progress}
    \begin{tikzpicture}
      \fill[dotoff] (0,0) rectangle (16,0.09);
      \ifnum\inserttotalframenumber>0
        \pgfmathsetmacro{\prog}{16*\insertframenumber/\inserttotalframenumber}
        \fill[accent] (0,0) rectangle (\prog,0.09);
      \fi
    \end{tikzpicture}
  \end{beamercolorbox}%
  \vskip 2pt
  \begin{beamercolorbox}[wd=\paperwidth, ht=2.4ex, dp=1.2ex, right]{footnum}
    \usebeamerfont{footline}\color{gray!60}\insertframenumber\,/\,\inserttotalframenumber\hspace{0.8em}
  \end{beamercolorbox}%
}

\setlength{\parindent}{0pt}
\setlength{\parskip}{6pt}

% ===================== CUSTOMIZAÇÃO DE BLOCOS =====================
\setbeamertemplate{blocks}[rounded][shadow=false]
\setbeamercolor{block title}{bg=dotoff, fg=accent}
\setbeamercolor{block body}{bg=bgcolor!90!white, fg=white!92!black}

\setbeamercolor{block title alerted}{bg=accent!60!black, fg=white}
\setbeamercolor{block body alerted}{bg=bgcolor!95!white, fg=white!92!black}

% ===================== CAIXA DE DESTAQUE =====================
% Uso: \notebox{texto da observação}
\newcommand{\notebox}[1]{%
  \begin{center}
  \begin{tabular}{@{}l@{\hspace{0.35cm}}p{0.82\textwidth}@{}}
    \textcolor{accent}{\rule{2pt}{0.55cm}} & \itshape\small\color{white!88!black} #1
  \end{tabular}
  \end{center}%
}

\usefonttheme[onlymath]{serif}
%
%      \title{Título da Aula}
%      \author
%
%      \begin{document}
%        ... seus \begin{frame} ... \end{frame} ...
%      \end{document}
%
%  IMPORTANTE: como este arquivo contém o \documentclass, ele
%  precisa ser a PRIMEIRA linha do .tex da aula (nenhum outro
%  comando pode vir antes do \input acima).
%
%  Compilar sempre a partir da pasta latex/ (para o caminho
%  relativo "templates/beamer.tex" ser encontrado):
%      cd latex && pdflatex aula01.tex
% ============================================================

% !TEX root = ../aula01.tex

\documentclass[11pt, aspectratio=169, xcolor=table]{beamer}

% ===================== PACOTES =====================
\usepackage[utf8]{inputenc}
\usepackage[T1]{fontenc}
\usepackage{lmodern}
\usepackage{newtxtext}
\usepackage{newtxmath}
\usepackage{tikz}
\usepackage{amsmath, amssymb}
\usepackage{booktabs}
\usepackage{bookmark}

% Atalho para "d" reto de diferenciais: dx, dy -> \dd x, \dd y
\newcommand{\dd}{\mathrm{d}}

% ===================== REMOVE ELEMENTOS PADRÃO DO BEAMER =====================
\setbeamertemplate{navigation symbols}{}

% ===================== CORES =====================
\definecolor{bgcolor}{RGB}{20, 20, 24}
\definecolor{accent}{RGB}{214, 165, 113}
\definecolor{dotoff}{RGB}{70, 70, 76}

\setbeamercolor{background canvas}{bg=bgcolor}
\setbeamercolor{normal text}{fg=white!92!black}
\setbeamercolor{title}{fg=white}
\setbeamercolor{frametitle}{fg=accent}
\setbeamercolor{footline}{fg=gray!60}
\setbeamercolor{itemize item}{fg=accent}
\setbeamercolor{itemize subitem}{fg=accent!70!white}
\setbeamercolor{progress}{bg=bgcolor}
\setbeamercolor{footnum}{bg=bgcolor}
\setbeamertemplate{itemize item}{\textbullet}

\setbeamerfont{frametitle}{series=\bfseries}
\setbeamerfont{footline}{size=\tiny}

\hypersetup{colorlinks=true, urlcolor=accent, linkcolor=accent}

% ===================== TÍTULO DO SLIDE =====================
\setbeamertemplate{frametitle}{
  \vspace{0.5cm}
  \hspace{0.4cm}\usebeamerfont{frametitle}\usebeamercolor[fg]{frametitle}\insertframetitle\\
  \vspace{0.12cm}
  \hspace{0.4cm}\textcolor{accent}{\rule{2.5cm}{0.6pt}}
}

% ===================== RODAPÉ =====================
\setbeamertemplate{footline}{%
  \begin{beamercolorbox}[wd=\paperwidth, ht=0.09cm, dp=0pt]{progress}
    \begin{tikzpicture}
      \fill[dotoff] (0,0) rectangle (16,0.09);
      \ifnum\inserttotalframenumber>0
        \pgfmathsetmacro{\prog}{16*\insertframenumber/\inserttotalframenumber}
        \fill[accent] (0,0) rectangle (\prog,0.09);
      \fi
    \end{tikzpicture}
  \end{beamercolorbox}%
  \vskip 2pt
  \begin{beamercolorbox}[wd=\paperwidth, ht=2.4ex, dp=1.2ex, right]{footnum}
    \usebeamerfont{footline}\color{gray!60}\insertframenumber\,/\,\inserttotalframenumber\hspace{0.8em}
  \end{beamercolorbox}%
}

\setlength{\parindent}{0pt}
\setlength{\parskip}{6pt}

% ===================== CUSTOMIZAÇÃO DE BLOCOS =====================
\setbeamertemplate{blocks}[rounded][shadow=false]
\setbeamercolor{block title}{bg=dotoff, fg=accent}
\setbeamercolor{block body}{bg=bgcolor!90!white, fg=white!92!black}

\setbeamercolor{block title alerted}{bg=accent!60!black, fg=white}
\setbeamercolor{block body alerted}{bg=bgcolor!95!white, fg=white!92!black}

% ===================== CAIXA DE DESTAQUE =====================
% Uso: \notebox{texto da observação}
\newcommand{\notebox}[1]{%
  \begin{center}
  \begin{tabular}{@{}l@{\hspace{0.35cm}}p{0.82\textwidth}@{}}
    \textcolor{accent}{\rule{2pt}{0.55cm}} & \itshape\small\color{white!88!black} #1
  \end{tabular}
  \end{center}%
}

\usefonttheme[onlymath]{serif}